\documentclass[12pt, a4paper]{article}
\usepackage[utf8]{inputenc}
\usepackage[T5]{fontenc}
\usepackage[vietnamese]{babel}
\usepackage{lmodern}
\usepackage{microtype}
\usepackage{amsmath, amssymb}
\usepackage{geometry}
\geometry{a4paper, top=1cm, bottom=1.2cm, left=1.2cm, right=1cm}
\usepackage{enumitem}
\usepackage{environ}
\usepackage{fancyhdr}
\usepackage{lastpage}
\usepackage{fontawesome5}
\usepackage{xcolor}
\usepackage{multicol}
\usepackage{tasks}
\AtBeginDocument{\fontsize{13pt}{15.6pt}\selectfont}
% In đậm phương án lựa chọn A, B, C, D
\settasks{label=\textbf{\Alph*.}, label-width=1.5em, item-indent=2em, column-sep=1em, label-offset=0.5em}

% TikZ Settings
\usepackage{tikz}
\usepackage{eso-pic}
\usepackage{tikz-3dplot}
\usepackage{pgfplots}
\pgfplotsset{compat=1.18}
\usetikzlibrary{3d, calc, intersections, arrows.meta, positioning, shapes.geometric, patterns, decorations.text, decorations.pathmorphing}

% Color Palette
\definecolor{primary}{HTML}{005F73}
\definecolor{secondary}{HTML}{0A9396}
\definecolor{accent}{HTML}{AE2012}
\definecolor{inquiry}{HTML}{6A4C93}
\definecolor{realworld}{HTML}{E76F51}
\definecolor{bg}{HTML}{F8F9FA}
\definecolor{solutionbg}{HTML}{E9F5F5}

\definecolor{LogoTeal}{RGB}{0, 102, 102}
\definecolor{LogoGold}{RGB}{218, 165, 32}
\definecolor{LogoDark}{RGB}{32, 50, 64}
\newcommand{\FourPointStar}[3]{
    \begin{scope}[shift={(#1,#2)}, scale=#3]
        \fill[LogoGold] (0,1) .. controls (0,0.1) and (0.1,0) .. (1,0)
                               .. controls (0.1,0) and (0,-0.1) .. (0,-1)
                               .. controls (0,-0.1) and (-0.1,0) .. (-1,0)
                               .. controls (-0.1,0) and (0,0.1) .. (0,1);
        \fill[white] (0,0) circle (0.12);
    \end{scope}
}

\newcommand{\PageLogoWatermark}{%
\begin{tikzpicture}[remember picture, overlay]
    \begin{scope}[shift={([yshift=-1.1cm]current page.center)}, opacity=0.15]
        \draw[LogoGold, line width=1.5pt] (0, 1.1) circle (2.8cm);
        \draw[LogoTeal, line width=1pt, decorate,
              decoration={snake, amplitude=0.15mm, segment length=1.5mm}]
            (0, 1.1) circle (2.93cm);
        \draw[LogoDark, line width=0.5pt] (0, 1.1) circle (3.06cm);
        \path[decorate,
              decoration={
                  text along path,
                  text={|\sffamily\bfseries\fontsize{10.8pt}{12pt}\selectfont\color{LogoDark}|Mathematics Since 2006},
                  text align=center,
                  text format delimiters={|}{}
              }]
            (0, -1.45) arc[start angle=-90, end angle=90, radius=2.55cm];
        \fill[LogoDark] (-1.6, -0.4) rectangle (-0.2, -0.2);
        \fill[LogoDark] (-1.6, -0.2) rectangle (-1.4, 0.8);
        \fill[LogoTeal] (-0.7, -0.2) rectangle (1.1, 0.0);
        \fill[LogoTeal] (0.1, -0.2) rectangle (0.3, -1.1);
        \fill[LogoGold]
            (-1.2, 0.4) .. controls (-1.05, 1.2) and (0.35, 0.7) .. (0.3, 1.5)
                        .. controls (0.4, 2.3) and (-0.8, 2.0) .. (-0.6, 3.2)
                        .. controls (-0.54, 3.42) and (-0.42, 3.52) .. (-0.18, 3.52)
                        .. controls (0.06, 3.52) and (0.08, 3.56) .. (0.2, 3.4)
                        .. controls (0.4, 3.1) and (-0.3, 2.2) .. (-0.2, 1.6)
                        .. controls (0.05, 0.8) and (-1.65, 1.2) .. (-1.2, 0.4);
        \FourPointStar{0.42}{3.48}{0.4}
        \FourPointStar{-1.712}{1.038}{0.4}
        \FourPointStar{0.441}{0.621}{0.8}
        \FourPointStar{1.15}{1.65}{0.6}
    \end{scope}
\end{tikzpicture}%
}
\AddToShipoutPictureBG{\PageLogoWatermark}

\usepackage[most]{tcolorbox}
\usepackage{etoolbox}

% --- MACRO TẠO TIÊU ĐỀ CÂU HỎI IN ĐẬM ---
\newcommand{\cau}[1]{\vspace{0.12cm}\noindent\textbf{Câu #1. }}

% --- MACRO PHÂN CHIA PHẦN THI ---
\newcommand{\PhanI}{
\vspace{0.25cm}
\noindent\textbf{\large\sffamily PHẦN I. CÂU TRẮC NGHIỆM NHIỀU PHƯƠNG ÁN LỰA CHỌN.}
\textit{Thí sinh trả lời từ câu 1 đến câu 16. Mỗi câu chọn một phương án đúng.}\par\vspace{0.2cm}
}
\newcommand{\PhanII}{
\vspace{0.3cm}
\noindent\textbf{\large\sffamily PHẦN II. CÂU TRẮC NGHIỆM ĐÚNG SAI.}
\textit{Thí sinh trả lời từ câu 17 đến câu 19. Trong mỗi ý a), b), c), d) ở mỗi câu, thí sinh chọn đúng hoặc sai.}\par\vspace{0.2cm}
}
\newcommand{\PhanIII}{
\vspace{0.3cm}
\noindent\textbf{\large\sffamily PHẦN III. CÂU TỰ LUẬN.}
\textit{Thí sinh trình bày chi tiết các bước giải vào giấy làm bài.}\par\vspace{0.2cm}
}

% --- LOGIC LỜI GIẢI VÀ DÒNG KẺ ---
\newbool{showsolutions}
\setbool{showsolutions}{true}

\newcommand{\filllines}[1]{}

\ifbool{showsolutions}{
\newtcolorbox{SolutionBox}{
enhanced, breakable, colback=solutionbg, colframe=secondary!50!black,
fonttitle=\bfseries\sffamily,
title={\faChalkboardTeacher\ Lời giải \& Chấm điểm},
borderline west={3pt}{0pt}{secondary!50!black}, frame hidden,
attach boxed title to top right={xshift=-0.5cm, yshift=-2mm},
boxed title style={colback=secondary!50!black, sharp corners}
}
}{
\NewEnviron{SolutionBox}{}
}

% --- HEADER CHUẨN ĐỀ THI ---
\pagestyle{fancy}
\fancyhf{}
\renewcommand{\headrulewidth}{0pt}
\cfoot{\thepage}

\begin{document}

% --- TIÊU ĐỀ ĐỀ THI ---
\noindent
\begin{minipage}[t]{0.52\textwidth}
\centering
SỞ GIÁO DỤC VÀ ĐÀO TẠO AN GIANG\\
\textbf{TRƯỜNG THPT NGUYỄN VĂN THOẠI}\\
\makebox[\linewidth][c]{\rule{4cm}{0.5pt}}\vspace{2mm}
\textit{(Đề thi gồm 4 trang)}
\end{minipage}
\hfill
\begin{minipage}[t]{0.45\textwidth}
\centering
\textbf{ĐỀ KIỂM TRA GIỮA HỌC KÌ I}\\
\textbf{NĂM HỌC 2026 -- 2027}\\
\textbf{Môn: TOÁN 11}\\
\textit{Thời gian làm bài: 90 phút}\\
\makebox[\linewidth][c]{\rule{4cm}{0.5pt}}\vspace{2mm}
\textbf{Mã đề thi: 11.GK1.TK4}
\end{minipage}
\vspace{0.3cm}

\PhanI

\cau{1} Góc có số đo $-450^\circ$ khi đổi sang đơn vị radian bằng
\begin{tasks}(4)
\task $-\dfrac{3\pi}{2}$.
\task $-\dfrac{7\pi}{3}$.
\task $-\dfrac{5\pi}{2}$.
\task $-\dfrac{5\pi}{4}$.
\end{tasks}

\cau{2} Cho góc lượng giác $x$ thỏa mãn $\cos x = \dfrac{3}{5}$ với $\dfrac{3\pi}{2} < x < 2\pi$. Giá trị của $\sin x$ bằng
\begin{tasks}(4)
\task $\dfrac{4}{5}$.
\task $-\dfrac{4}{5}$.
\task $-\dfrac{3}{5}$.
\task $\dfrac{16}{25}$.
\end{tasks}

\cau{3} Khẳng định nào sau đây là khẳng định \textbf{đúng} với mọi góc $a$?
\begin{tasks}(2)
\task $\cos 2a = \cos^2 a + \sin^2 a$.
\task $\sin 2a = \sin a\cos a$.
\task $\cos 2a = 1 - 2\cos^2 a$.
\task $\cos 2a = 2\cos^2 a - 1$.
\end{tasks}

\cau{4} Tập xác định của hàm số $y = \tan x$ là
\begin{tasks}(2)
\task $D = \mathbb{R}\setminus\left\{\dfrac{\pi}{2} + k\pi,\, k\in\mathbb{Z}\right\}$.
\task $D = \mathbb{R}$.
\task $D = \mathbb{R}\setminus\{k\pi,\, k\in\mathbb{Z}\}$.
\task $D = [-1; 1]$.
\end{tasks}

\cau{5} Nghiệm của phương trình lượng giác $\cos x = -1$ là
\begin{tasks}(2)
\task $x = k2\pi,\, k\in\mathbb{Z}$.
\task $x = \pi + k2\pi,\, k\in\mathbb{Z}$.
\task $x = -\dfrac{\pi}{2} + k2\pi,\, k\in\mathbb{Z}$.
\task $x = \pi + k\pi,\, k\in\mathbb{Z}$.
\end{tasks}

\cau{6} Phương trình lượng giác $2\cos x + 1 = 0$ có tập nghiệm là
\begin{tasks}(2)
\task $S = \left\{\pm\dfrac{\pi}{3} + k2\pi,\, k\in\mathbb{Z}\right\}$.
\task $S = \left\{\pm\dfrac{5\pi}{6} + k2\pi,\, k\in\mathbb{Z}\right\}$.
\task $S = \left\{\pm\dfrac{2\pi}{3} + k2\pi,\, k\in\mathbb{Z}\right\}$.
\task $S = \left\{\pm\dfrac{\pi}{6} + k2\pi,\, k\in\mathbb{Z}\right\}$.
\end{tasks}

\noindent
\begin{minipage}[t]{0.70\linewidth}
\cau{7} Hội Lim (tỉnh Bắc Ninh) tổ chức vào mùa xuân thường có trò chơi đánh đu. Khi người chơi đu nhún đều, cây đu sẽ đưa người chơi đu dao động quanh vị trí cân bằng (hình vẽ bên). Nghiên cứu trò chơi này, người ta thấy khoảng cách $h$ (đơn vị: mét) từ vị trí người chơi đu đến vị trí cân bằng được biểu diễn qua thời gian $t$ (giây, với $t \ge 0$) bởi hệ thức $h = |d|$ với $d = 4\cos\left[\dfrac{\pi}{4}(2t - 1)\right]$, trong đó ta quy ước $d > 0$ khi vị trí cân bằng ở phía sau lưng người chơi đu và $d < 0$ trong trường hợp ngược lại. Tìm các thời điểm $t$ theo tham số $k \in \mathbb{Z}$ ($k \ge 0$) để khoảng cách $h = 4\text{ m}$.
\begin{tasks}(2)
\task $t = \dfrac{4k-1}{2},\, k\in\mathbb{Z}, k \ge 1$.
\task $t = \dfrac{4k+1}{2},\, k\in\mathbb{Z}, k \ge 0$.
\task $t = 2k + 1,\, k\in\mathbb{Z}, k \ge 0$.
\task $t = \dfrac{2k+1}{4},\, k\in\mathbb{Z}, k \ge 0$.
\end{tasks}
\end{minipage}\hfill
\begin{minipage}[t]{0.28\linewidth}
\centering
\begin{tikzpicture}[scale=0.82, line join=round, line cap=round]
  % Gia treo phia tren
  \coordinate (O) at (0, 3.8);
  
  % Cung dao dong cua cay du
  \draw[line width=1.3pt, secondary!85!black] (-2.2, 0.72) arc (-125:-55:3.8);
  \fill[LogoDark] (-2.2, 0.72) circle (2.2pt);
  \fill[LogoDark] (2.2, 0.72) circle (2.2pt);
  
  % Duong thang dung vi tri can bang
  \draw[dashed, line width=0.8pt, gray!75!black] (0, 4.1) -- (0, -0.05);
  \node[right, font=\scriptsize\itshape, text=primary!90!black, rotate=90] at (0.24, 1.8) {Vị trí cân bằng};
  
  % Dinh treo O
  \fill[primary] (O) circle (2.5pt);
  \draw[thin, gray!50, dashed] (O) -- (1.35, 0.46);
  
  % Day du dang lech ve phia truoc
  \coordinate (P) at (-1.35, 0.46);
  \draw[line width=1.5pt, primary!90!black] (O) -- (P);
  
  % Nguoi choi du & ban du
  \draw[line width=2.4pt, brown!80!black] ($(P)+(-0.35,-0.05)$) -- ($(P)+(0.35,0.05)$);
  \fill[accent] ($(P)+(0,0.46)$) circle (3pt); % Dau
  \draw[line width=1.6pt, accent] ($(P)+(0,0.12)$) -- ($(P)+(0,0.42)$); % Than
  \draw[line width=1pt, accent] ($(P)+(-0.16,0.30)$) -- ($(P)+(0,0.36)$) -- ($(P)+(0.12,0.52)$); % Tay
  \draw[line width=1.3pt, primary!80!black] ($(P)+(-0.1,0.0)$) -- ($(P)+(0,0.14)$) -- ($(P)+(0.1,0.0)$); % Chan
  
  % Mui ten khoang cach h
  \draw[<->, >=stealth, line width=0.9pt, accent] (-1.35, 0.15) -- (0, 0.15);
  \node[below, font=\small\bfseries, text=accent] at (-0.68, 0.15) {$h$};
\end{tikzpicture}
\end{minipage}

\cau{8} Cho dãy số $(u_n)$ có số hạng tổng quát $u_n = \dfrac{3n + 1}{n + 2}$. Giá trị của số hạng thứ ba $u_3$ bằng
\begin{tasks}(4)
\task $\dfrac{7}{5}$.
\task $2$.
\task $\dfrac{8}{5}$.
\task $\dfrac{5}{2}$.
\end{tasks}

\cau{9} Dãy số nào sau đây là một dãy số tăng?
\begin{tasks}(2)
\task $15;\, 12;\, 9;\, 6;\, 3;\,\dots$
\task $-2;\,-5;\,-8;\,-11;\,-14;\,\dots$
\task $3;\,-3;\, 3;\,-3;\, 3;\,\dots$
\task $2;\, 5;\, 8;\, 11;\, 14;\,\dots$
\end{tasks}

\cau{10} Cho cấp số cộng $(u_n)\colon 3;\, 8;\, 13;\, 18;\, 23;\,\dots$. Công sai $d$ của cấp số cộng đã cho bằng
\begin{tasks}(4)
\task $d = 5$.
\task $d = -5$.
\task $d = 3$.
\task $d = \dfrac{8}{3}$.
\end{tasks}

\cau{11} Cho cấp số nhân $(u_n)$ có số hạng đầu $u_1 = 9$ và công bội $q = -\dfrac{1}{3}$. Số hạng thứ hai $u_2$ bằng
\begin{tasks}(4)
\task $u_2 = 3$.
\task $u_2 = -27$.
\task $u_2 = -3$.
\task $u_2 = \dfrac{26}{3}$.
\end{tasks}

\cau{12} Cho cấp số cộng $(u_n)$ có số hạng đầu $u_1 = -3$ và công sai $d = 5$. Số $142$ là số hạng thứ mấy của cấp số cộng đã cho?
\begin{tasks}(4)
\task Thứ $29$.
\task Thứ $31$.
\task Thứ $28$.
\task Thứ $30$.
\end{tasks}

\noindent
\begin{minipage}[t]{0.70\linewidth}
\cau{13} Một người thợ khoan một cái giếng sâu $25\text{ m}$. Giá tiền của mét khoan đầu tiên là $120\,000\text{ đồng}$. Kể từ mét khoan thứ hai trở đi, giá của mỗi mét sau tăng thêm $30\,000\text{ đồng}$ so với mét ngay trước đó. Hỏi tổng chi phí để người đó khoan hết $25\text{ m}$ giếng là bao nhiêu tiền?
\begin{tasks}(2)
\task $12\text{ triệu đồng}$.
\task $11{,}5\text{ triệu đồng}$.
\task $12{,}5\text{ triệu đồng}$.
\task $10{,}8\text{ triệu đồng}$.
\end{tasks}
\end{minipage}\hfill
\begin{minipage}[t]{0.28\linewidth}
\centering
\begin{tikzpicture}[scale=0.68, line join=round, line cap=round]
  % Mat dat
  \fill[green!50!brown!20] (-2.2, 0) rectangle (2.2, -0.15);
  \draw[line width=1.1pt, green!60!black] (-2.2, 0) -- (2.2, 0);
  \node[above, font=\tiny\itshape, text=green!50!black] at (-1.3, 0.02) {Mặt đất};
  
  % Thap khoan mat dat
  \draw[line width=1.1pt, gray!80!black] (-0.6, 0) -- (0, 1.4) -- (0.6, 0);
  \draw[line width=0.7pt, gray!70!black] (-0.4, 0.5) -- (0.4, 0.5);
  \draw[line width=0.7pt, gray!70!black] (-0.22, 0.95) -- (0.22, 0.95);
  \fill[primary] (0, 1.4) circle (1.8pt);
  \draw[line width=1.2pt, primary] (0, 1.4) -- (0, 0);
  
  % Long gieng khoan
  \fill[cyan!8] (-0.45, 0) rectangle (0.45, -2.6);
  \draw[line width=1pt, LogoDark] (-0.45, 0) -- (-0.45, -2.6);
  \draw[line width=1pt, LogoDark] (0.45, 0) -- (0.45, -2.6);
  \draw[line width=1pt, LogoDark] (-0.45, -2.6) arc (180:360:0.45 and 0.15);
  
  % Truc mui khoan
  \draw[line width=1.2pt, primary] (0, 0) -- (0, -2.4);
  \fill[accent] (-0.18, -2.4) -- (0.18, -2.4) -- (0, -2.7) -- cycle;
  
  % Cac vach phan tang
  \draw[dashed, thin, gray!60] (-0.45, -0.5) -- (0.45, -0.5);
  \draw[dashed, thin, gray!60] (-0.45, -1.0) -- (0.45, -1.0);
  \node[font=\tiny, text=gray!80!black] at (0, -1.5) {$\vdots$};
  \draw[dashed, thin, gray!60] (-0.45, -2.1) -- (0.45, -2.1);
  
  % Chu thich
  \draw[->, >=stealth, thin, primary] (0.9, -0.25) -- (0.48, -0.25);
  \node[right, font=\tiny] at (0.85, -0.25) {$120\text{k}$};
  \draw[->, >=stealth, thin, primary] (0.9, -0.75) -- (0.48, -0.75);
  \node[right, font=\tiny] at (0.85, -0.75) {$150\text{k}$};
  \draw[->, >=stealth, thin, primary] (0.9, -2.35) -- (0.48, -2.35);
  \node[right, font=\tiny] at (0.85, -2.35) {$u_{25}$};
  
  % Do sau 25m
  \draw[<->, >=stealth, line width=0.8pt, accent] (-0.8, 0) -- (-0.8, -2.6);
  \node[left, font=\scriptsize\bfseries, text=accent] at (-0.8, -1.3) {$25\text{m}$};
\end{tikzpicture}
\end{minipage}

\cau{14} Hình nào sau đây là hình biểu diễn đúng của một hình chóp tứ giác?
\begin{tasks}(4)
\task \begin{tikzpicture}[scale=0.48, baseline=(current bounding box.center), line join=round, line cap=round]
\coordinate (A) at (0,0);
\coordinate (B) at (1.2,-0.8);
\coordinate (C) at (3.2,-0.6);
\coordinate (D) at (2.0,0.3);
\coordinate (S) at (1.4,2.5);
\draw[thick] (S)--(A)--(B)--(C)--(D)--(S)--(B) (S)--(C) (A)--(B)--(C)--(D)--cycle;
\end{tikzpicture}
\task \begin{tikzpicture}[scale=0.48, baseline=(current bounding box.center), line join=round, line cap=round]
\coordinate (A) at (0,0);
\coordinate (B) at (1.2,-0.8);
\coordinate (C) at (3.2,-0.6);
\coordinate (D) at (2.0,0.3);
\coordinate (S) at (1.4,2.5);
\draw[dashed, thin] (A)--(D)--(C) (S)--(D);
\draw[thick] (S)--(A)--(B)--(C)--(S)--(B);
\end{tikzpicture}
\task \begin{tikzpicture}[scale=0.48, baseline=(current bounding box.center), line join=round, line cap=round]
\coordinate (A) at (0,0);
\coordinate (B) at (1.1,-0.8);
\coordinate (C) at (2.9,0);
\coordinate (S) at (1.3,2.4);
\draw[dashed, thin] (A)--(C);
\draw[thick] (S)--(A)--(B)--(C)--(S)--(B);
\end{tikzpicture}
\task \begin{tikzpicture}[scale=0.48, baseline=(current bounding box.center), line join=round, line cap=round]
\coordinate (A) at (0,0);
\coordinate (B) at (1.4,-0.8);
\coordinate (C) at (3.2,0);
\coordinate (D) at (1.8,0.8);
\draw[thick] (A)--(B)--(C)--(D)--cycle (A)--(C) (B)--(D);
\end{tikzpicture}
\end{tasks}

\cau{15} Cho hình chóp $S.ABCD$ có đáy $ABCD$ là hình thang với đáy lớn $AD$, đáy nhỏ $BC$ ($AD \parallel BC$). Hai cạnh bên $AB$ và $CD$ kéo dài cắt nhau tại điểm $E$. Giao tuyến của hai mặt phẳng $(SAB)$ và $(SCD)$ là đường thẳng nào sau đây?
\begin{tasks}(4)
\task $SO$.
\task $AC$.
\task $SE$.
\task $BD$.
\end{tasks}

\cau{16} Cho hình chóp $S.ABCD$ có đáy $ABCD$ là hình thang với đáy lớn $AB$, đáy nhỏ $CD$ ($AB \parallel CD$). Gọi $O$ là giao điểm của hai đường chéo $AC$ và $BD$. Trên cạnh bên $SC$ lấy điểm $M$ ($M$ không trùng với $S$ và $C$). Trong mặt phẳng $(SAC)$, đường thẳng $SO$ cắt đường thẳng $AM$ tại điểm $P$. Giao điểm của đường thẳng $AM$ với mặt phẳng $(SBD)$ là
\begin{tasks}(4)
\task Điểm $O$.
\task Điểm $M$.
\task Điểm $S$.
\task Điểm $P$.
\end{tasks}

\PhanII

\cau{17} Cho hàm số lượng giác $y = f(x) = \tan x$. Xét tính đúng sai của các mệnh đề sau:
\begin{enumerate}[label=\alph*),leftmargin=*]
\item Tập xác định của hàm số là $D = \mathbb{R}\setminus\left\{\dfrac{\pi}{2} + k\pi \;\middle|\; k\in\mathbb{Z}\right\}$.
\item Giá trị của hàm số tại điểm $x_0 = \dfrac{\pi}{4}$ bằng $1$.
\item Chu kì tuần hoàn của hàm số là $T = 2\pi$.
\item Trên khoảng $\left(-\dfrac{\pi}{2}; \dfrac{\pi}{2}\right)$, phương trình $f(x) = \sqrt{3}$ có đúng $1$ nghiệm phân biệt là $x = \dfrac{\pi}{3}$.
\end{enumerate}

\cau{18} Để đón đầu cuộc cách mạng công nghiệp bán dẫn, một tập đoàn viễn thông quyết định rót vốn đầu tư dài hạn xây dựng trung tâm nghiên cứu và phát triển vi mạch công nghệ cao. Bắt đầu từ quý I năm đầu tiên, tập đoàn tiến hành giải ngân số vốn ban đầu là $16\text{ tỷ đồng}$. Theo kế hoạch đầu tư được phê duyệt, để bám sát tiến độ nghiên cứu, số vốn giải ngân ở mỗi quý tiếp theo luôn gấp $1{,}5$ lần số vốn giải ngân của quý liền trước đó. Gọi $u_n$ (đơn vị: tỷ đồng) là số tiền giải ngân ở quý thứ $n$ ($n \in \mathbb{N}^*$), và $S_n = u_1 + u_2 + \dots + u_n$ là tổng số vốn lũy kế đã giải ngân sau $n$ quý đầu tiên.
\begin{center}
\begin{tikzpicture}[scale=0.88, line join=round, line cap=round]
  % Khung man hinh may tinh
  \filldraw[fill=bg, draw=primary!70!black, line width=1.2pt, rounded corners=4pt] (-0.8,-0.7) rectangle (9.6,4.0);
  \node[font=\sffamily\bfseries\scriptsize, text=primary] at (4.4,3.6) {KẾ HOẠCH GIẢI NGÂN VỐN ĐẦU TƯ R\&D THEO TỪNG QUÝ (TỶ ĐỒNG)};
  
  % Luoi truc toa do
  \draw[->, >=stealth, thick, gray!80!black] (0,0) -- (9.0,0) node[below, font=\scriptsize] {Quý};
  \draw[->, >=stealth, thick, gray!80!black] (0,0) -- (0,3.3) node[left, font=\scriptsize] {Tỷ đồng};
  
  % Cac moc tren truc hoanh
  \foreach \x/\q in {1.2/Quý 1, 2.6/Quý 2, 4.0/Quý 3, 5.4/Quý 4, 6.8/Quý 5, 8.2/Quý 6} {
    \draw[gray!60!black] (\x,0.08) -- (\x,-0.08);
    \node[below, font=\scriptsize] at (\x,-0.1) {\q};
  }
  % Cac moc tren truc tung
  \foreach \y/\val in {0.4/16, 0.8/24, 1.2/36, 1.8/54, 2.7/81} {
    \draw[dashed, thin, gray!30] (0,\y) -- (8.5,\y);
    \draw[gray!60!black] (0.08,\y) -- (-0.08,\y);
    \node[left, font=\scriptsize] at (-0.08,\y) {\val};
  }
  
  % Ve cac cot gia tri
  \filldraw[fill=secondary!30, draw=secondary!80!black] (0.9,0) rectangle (1.5,0.4);
  \node[above, font=\scriptsize\bfseries, text=secondary!90!black] at (1.2,0.4) {16};
  
  \filldraw[fill=secondary!30, draw=secondary!80!black] (2.3,0) rectangle (2.9,0.8);
  \node[above, font=\scriptsize\bfseries, text=secondary!90!black] at (2.6,0.8) {24};
  
  \filldraw[fill=secondary!30, draw=secondary!80!black] (3.7,0) rectangle (4.3,1.2);
  \node[above, font=\scriptsize\bfseries, text=secondary!90!black] at (4.0,1.2) {36};
  
  \filldraw[fill=secondary!30, draw=secondary!80!black] (5.1,0) rectangle (5.7,1.8);
  \node[above, font=\scriptsize\bfseries, text=secondary!90!black] at (5.4,1.8) {54};
  
  \filldraw[fill=secondary!30, draw=secondary!80!black] (6.5,0) rectangle (7.1,2.7);
  \node[above, font=\scriptsize\bfseries, text=secondary!90!black] at (6.8,2.7) {81};
  
  % Duong cong xu huong tang truong cap so nhan
  \draw[domain=1.2:8.2, smooth, variable=\t, line width=1.5pt, accent] 
    plot ({\t}, {0.4 * pow(1.5, (\t-1.2)/1.4)});
  \fill[accent] (1.2,0.4) circle (2.2pt);
  \fill[accent] (6.8,2.7) circle (2.2pt);
  \node[above left, font=\scriptsize\bfseries, text=accent] at (8.5,3.1) {$u_n = 16 \cdot (1{,}5)^{n-1}$};
\end{tikzpicture}
\end{center}
Xét tính đúng sai của các khẳng định sau:
\begin{enumerate}[label=\alph*),leftmargin=*]
\item Dãy số $(u_n)$ là một cấp số nhân có số hạng đầu $u_1 = 16$ tỷ đồng và công bội $q = 1{,}5$.
\item Công thức tính số tiền giải ngân ở quý thứ $n$ là $u_n = 16 \cdot (1{,}5)^n$ tỷ đồng với mọi $n \ge 1$.
\item Số tiền giải ngân của riêng quý thứ tư là $54$ tỷ đồng, và tổng số vốn tập đoàn đã giải ngân tích lũy sau $1$ năm đầu tiên (gồm 4 quý) là $130$ tỷ đồng.
\item Giả sử tổng ngân sách dự kiến của cả giai đoạn nghiên cứu là $500$ tỷ đồng. Sau đúng $6$ quý giải ngân, tổng số vốn giải ngân lũy kế của tập đoàn đã vượt mức ngân sách dự kiến này.
\end{enumerate}

\cau{19} Cho hình chóp $S.ABCD$ có đáy $ABCD$ là hình bình hành tâm $O$. Gọi $M$ là trung điểm của cạnh $CD$. Gọi $G$ và $E$ lần lượt là trọng tâm của tam giác $SCD$ và tam giác $BCD$ (tham khảo hình vẽ bên dưới).
\begin{center}
\begin{tikzpicture}[scale=0.9, line join=round, line cap=round]
  % Dinh hinh binh hanh va dinh chóp
  \coordinate (A) at (0,0);
  \coordinate (B) at (4.6,-0.6);
  \coordinate (D) at (1.6,1.8);
  \coordinate (C) at ($(B)+(D)-(A)$);
  \coordinate (S) at (2.2,4.8);
  
  % Tam O va trung diem M
  \coordinate (O) at (intersection of A--C and B--D);
  \coordinate (M) at ($(C)!0.5!(D)$);
  
  % Cac trong tam G tren SM va E tren BM
  \coordinate (G) at ($(S)!0.667!(M)$);
  \coordinate (E) at ($(B)!0.667!(M)$);
  
  % Net dut khuat
  \draw[dashed, thick] (A)--(D)--(C) (S)--(D) (A)--(C) (B)--(D) (S)--(O);
  \draw[dashed, thick, primary] (S)--(M) (B)--(M) (O)--(M);
  \draw[dashed, thick, accent] (G)--(E);
  
  % Net lien nhin thay
  \draw[thick] (S)--(A)--(B)--(C)--(S)--(B);
  
  % Cac diem
  \fill (S) circle (2pt) node[above] {$S$};
  \fill (A) circle (2pt) node[below left] {$A$};
  \fill (B) circle (2pt) node[below right] {$B$};
  \fill (C) circle (2pt) node[right] {$C$};
  \fill (D) circle (2pt) node[above left] {$D$};
  \fill (O) circle (2pt) node[below] {$O$};
  \fill[primary] (M) circle (2pt) node[above right] {$M$};
  \fill[accent] (G) circle (2.2pt) node[left] {$G$};
  \fill[accent] (E) circle (2.2pt) node[below] {$E$};
\end{tikzpicture}
\end{center}
Xét tính đúng sai của các khẳng định sau:
\begin{enumerate}[label=\alph*),leftmargin=*]
\item Điểm $G$ thuộc mặt phẳng $(SCD)$.
\item Giao tuyến của hai mặt phẳng $(SCD)$ và $(SBD)$ là đường thẳng $SD$.
\item Đường thẳng $OM$ song song với đường thẳng $BC$.
\item Đường thẳng $GE$ cắt mặt phẳng $(SAB)$ tại điểm $P$ thì $P \in SB$.
\end{enumerate}

\PhanIII

\cau{20} Giải phương trình lượng giác:
$$\cos\left(2x + \dfrac{\pi}{4}\right) = -\dfrac{\sqrt{2}}{2}$$
\filllines{5}

\cau{21} Cho hình chóp $S.ABCD$ có đáy $ABCD$ là hình bình hành tâm $O$. Gọi $M$ là trung điểm của cạnh bên $SC$. Xác định giao điểm $K$ của đường thẳng $SD$ với mặt phẳng $(ABM)$.
\filllines{6}

\clearpage
\section*{BẢNG ĐÁP ÁN}
\noindent\textbf{Phần I. Trắc nghiệm nhiều phương án lựa chọn}
\begin{center}
\small
\renewcommand{\arraystretch}{1.2}
\begin{tabular}{c|*{16}{c|}}
Câu & 1 & 2 & 3 & 4 & 5 & 6 & 7 & 8 & 9 & 10 & 11 & 12 & 13 & 14 & 15 & 16 \\\hline
Đáp án & C & B & D & A & B & C & B & B & D & A & C & D & A & B & C & D \\
\end{tabular}
\end{center}
\noindent\textbf{Phần II. Trắc nghiệm đúng sai}
\begin{center}
\small
\renewcommand{\arraystretch}{1.2}
\begin{tabular}{c|c|c|c|c}
Câu & a) & b) & c) & d) \\\hline
17 & Đ & Đ & S & Đ \\
18 & Đ & S & Đ & S \\
19 & Đ & Đ & Đ & S \\
\end{tabular}
\end{center}

\section*{LỜI GIẢI CHI TIẾT}

\noindent\textbf{Câu 1. Lời giải chi tiết}\par
\begin{tcolorbox}[enhanced,breakable,colback=solutionbg,colframe=secondary!50!black,frame hidden,borderline west={3pt}{0pt}{secondary!50!black}]
Áp dụng công thức đổi số đo góc từ độ sang radian: $\alpha_{\text{rad}} = a^\circ \cdot \dfrac{\pi}{180^\circ}$.\\
Ta có:
$$-450^\circ = -450 \cdot \dfrac{\pi}{180} = -\dfrac{5\pi}{2}\text{ (rad)}.$$
\textbf{Chọn C.}
\end{tcolorbox}

\noindent\textbf{Câu 2. Lời giải chi tiết}\par
\begin{tcolorbox}[enhanced,breakable,colback=solutionbg,colframe=secondary!50!black,frame hidden,borderline west={3pt}{0pt}{secondary!50!black}]
Áp dụng hệ thức cơ bản giữa $\sin x$ và $\cos x$:
$$\sin^2 x + \cos^2 x = 1 \Rightarrow \sin^2 x = 1 - \cos^2 x = 1 - \left(\dfrac{3}{5}\right)^2 = 1 - \dfrac{9}{25} = \dfrac{16}{25}.$$
Vì $\dfrac{3\pi}{2} < x < 2\pi$ (góc lượng giác $x$ thuộc góc phần tư thứ IV) nên $\sin x < 0$. Do đó:
$$\sin x = -\sqrt{\dfrac{16}{25}} = -\dfrac{4}{5}.$$
\textbf{Chọn B.}
\end{tcolorbox}

\noindent\textbf{Câu 3. Lời giải chi tiết}\par
\begin{tcolorbox}[enhanced,breakable,colback=solutionbg,colframe=secondary!50!black,frame hidden,borderline west={3pt}{0pt}{secondary!50!black}]
Theo công thức nhân đôi lượng giác đối với côsin:
$$\cos 2a = \cos^2 a - \sin^2 a = 2\cos^2 a - 1 = 1 - 2\sin^2 a.$$
Do đó khẳng định ở phương án D là khẳng định đúng.\\
Các khẳng định còn lại đều sai vì: $\cos^2 a + \sin^2 a = 1$; $\sin 2a = 2\sin a\cos a$; $1 - 2\cos^2 a = -\cos 2a$.\\
\textbf{Chọn D.}
\end{tcolorbox}

\noindent\textbf{Câu 4. Lời giải chi tiết}\par
\begin{tcolorbox}[enhanced,breakable,colback=solutionbg,colframe=secondary!50!black,frame hidden,borderline west={3pt}{0pt}{secondary!50!black}]
Hàm số lượng giác $y = \tan x = \dfrac{\sin x}{\cos x}$ xác định khi và chỉ khi:
$$\cos x \ne 0 \Leftrightarrow x \ne \dfrac{\pi}{2} + k\pi\quad (k\in\mathbb{Z}).$$
Vậy tập xác định của hàm số là $D = \mathbb{R}\setminus\left\{\dfrac{\pi}{2} + k\pi,\, k\in\mathbb{Z}\right\}$.\\
\textbf{Chọn A.}
\end{tcolorbox}

\noindent\textbf{Câu 5. Lời giải chi tiết}\par
\begin{tcolorbox}[enhanced,breakable,colback=solutionbg,colframe=secondary!50!black,frame hidden,borderline west={3pt}{0pt}{secondary!50!black}]
Theo công thức nghiệm của phương trình lượng giác cơ bản trường hợp đặc biệt:
$$\cos x = -1 \Leftrightarrow x = \pi + k2\pi\quad (k\in\mathbb{Z}).$$
\textbf{Chọn B.}
\end{tcolorbox}

\noindent\textbf{Câu 6. Lời giải chi tiết}\par
\begin{tcolorbox}[enhanced,breakable,colback=solutionbg,colframe=secondary!50!black,frame hidden,borderline west={3pt}{0pt}{secondary!50!black}]
Ta có: $2\cos x + 1 = 0 \Leftrightarrow \cos x = -\dfrac{1}{2} = \cos\left(\dfrac{2\pi}{3}\right)$.
$$\Leftrightarrow x = \pm\dfrac{2\pi}{3} + k2\pi\quad (k\in\mathbb{Z}).$$
Vậy tập nghiệm của phương trình là $S = \left\{\pm\dfrac{2\pi}{3} + k2\pi,\, k\in\mathbb{Z}\right\}$.\\
\textbf{Chọn C.}
\end{tcolorbox}

\noindent\textbf{Câu 7. Lời giải chi tiết}\par
\begin{tcolorbox}[enhanced,breakable,colback=solutionbg,colframe=secondary!50!black,frame hidden,borderline west={3pt}{0pt}{secondary!50!black}]
Khoảng cách từ người chơi đu đến vị trí cân bằng là $h = |d|$.\\
Theo đề bài, ta cần tìm thời điểm $t$ để $h = 4\text{ m}$:
$$|d| = 4 \Leftrightarrow \left|4\cos\left[\dfrac{\pi}{4}(2t - 1)\right]\right| = 4 \Leftrightarrow \left|\cos\left[\dfrac{\pi}{4}(2t - 1)\right]\right| = 1.$$
Điều này tương đương với:
$$\cos\left[\dfrac{\pi}{4}(2t - 1)\right] = \pm 1 \Leftrightarrow \dfrac{\pi}{4}(2t - 1) = k\pi\quad (k\in\mathbb{Z}).$$
$$\Leftrightarrow 2t - 1 = 4k \Leftrightarrow 2t = 4k + 1 \Leftrightarrow t = \dfrac{4k + 1}{2}\quad (k\in\mathbb{Z}).$$
Vì thời gian $t \ge 0$ nên $\dfrac{4k+1}{2} \ge 0 \Leftrightarrow k \ge -\dfrac{1}{4}$. Do $k \in \mathbb{Z}$ nên $k \ge 0$ ($k \in \{0, 1, 2, \dots\}$).\\
Vậy các thời điểm thỏa mãn yêu cầu là $t = \dfrac{4k+1}{2}$ ($k \in \mathbb{Z}, k \ge 0$).\\
\textbf{Chọn B.}
\end{tcolorbox}

\noindent\textbf{Câu 8. Lời giải chi tiết}\par
\begin{tcolorbox}[enhanced,breakable,colback=solutionbg,colframe=secondary!50!black,frame hidden,borderline west={3pt}{0pt}{secondary!50!black}]
Thay $n = 3$ vào công thức số hạng tổng quát của dãy số:
$$u_3 = \dfrac{3\cdot 3 + 1}{3 + 2} = \dfrac{9 + 1}{5} = \dfrac{10}{5} = 2.$$
\textbf{Chọn B.}
\end{tcolorbox}

\noindent\textbf{Câu 9. Lời giải chi tiết}\par
\begin{tcolorbox}[enhanced,breakable,colback=solutionbg,colframe=secondary!50!black,frame hidden,borderline west={3pt}{0pt}{secondary!50!black}]
Dãy số $2;\, 5;\, 8;\, 11;\, 14;\,\dots$ có hiệu của hai số liên tiếp:
$$u_{n+1} - u_n = 3 > 0 \Leftrightarrow u_{n+1} > u_n\quad (\forall n \ge 1),$$
nên là dãy số tăng. Dãy A và B là các dãy số giảm; dãy C là dãy đan dấu (không tăng, không giảm).\\
\textbf{Chọn D.}
\end{tcolorbox}

\noindent\textbf{Câu 10. Lời giải chi tiết}\par
\begin{tcolorbox}[enhanced,breakable,colback=solutionbg,colframe=secondary!50!black,frame hidden,borderline west={3pt}{0pt}{secondary!50!black}]
Công sai của cấp số cộng bằng hiệu của số hạng thứ hai trừ đi số hạng đầu:
$$d = u_2 - u_1 = 8 - 3 = 5.$$
\textbf{Chọn A.}
\end{tcolorbox}

\noindent\textbf{Câu 11. Lời giải chi tiết}\par
\begin{tcolorbox}[enhanced,breakable,colback=solutionbg,colframe=secondary!50!black,frame hidden,borderline west={3pt}{0pt}{secondary!50!black}]
Theo định nghĩa cấp số nhân, số hạng thứ hai được tính bởi:
$$u_2 = u_1 \cdot q = 9 \cdot \left(-\dfrac{1}{3}\right) = -3.$$
\textbf{Chọn C.}
\end{tcolorbox}

\noindent\textbf{Câu 12. Lời giải chi tiết}\par
\begin{tcolorbox}[enhanced,breakable,colback=solutionbg,colframe=secondary!50!black,frame hidden,borderline west={3pt}{0pt}{secondary!50!black}]
Công thức số hạng tổng quát của cấp số cộng: $u_n = u_1 + (n - 1)d$. Do đó:
$$142 = -3 + (n - 1)\cdot 5 \Leftrightarrow 5(n - 1) = 145 \Leftrightarrow n - 1 = 29 \Leftrightarrow n = 30.$$
Vậy $142$ là số hạng thứ $30$ của cấp số cộng đã cho.\\
\textbf{Chọn D.}
\end{tcolorbox}

\noindent\textbf{Câu 13. Lời giải chi tiết}\par
\begin{tcolorbox}[enhanced,breakable,colback=solutionbg,colframe=secondary!50!black,frame hidden,borderline west={3pt}{0pt}{secondary!50!black}]
Giá tiền của từng mét khoan lập thành một cấp số cộng $(u_n)$ với:
\begin{itemize}[leftmargin=1.5em]
\item Số hạng đầu: $u_1 = 120\,000$ (đồng);
\item Công sai: $d = 30\,000$ (đồng);
\item Số số hạng: $n = 25$ (mét).
\end{itemize}
Tổng chi phí để người đó khoan hết $25\text{ m}$ giếng là tổng của $25$ số hạng đầu tiên:
$$S_{25} = \dfrac{25}{2}\left[2u_1 + (25 - 1)d\right] = \dfrac{25}{2}\left[2 \cdot 120\,000 + 24 \cdot 30\,000\right]$$
$$= \dfrac{25}{2}\left[240\,000 + 720\,000\right] = \dfrac{25}{2} \cdot 960\,000 = 25 \cdot 480\,000 = 12\,000\,000\text{ (đồng)}.$$
Đổi $12\,000\,000\text{ đồng} = 12\text{ triệu đồng}$.\\
\textbf{Chọn A.}
\end{tcolorbox}

\noindent\textbf{Câu 14. Lời giải chi tiết}\par
\begin{tcolorbox}[enhanced,breakable,colback=solutionbg,colframe=secondary!50!black,frame hidden,borderline west={3pt}{0pt}{secondary!50!black}]
Hình biểu diễn của một hình chóp tứ giác là hình chóp có đáy là tứ giác lồi, gồm $1$ đỉnh và đáy $4$ đỉnh ($5$ đỉnh, $8$ cạnh). Khi biểu diễn trong không gian:
\begin{itemize}[leftmargin=1.5em]
\item Các đường nhìn thấy vẽ bằng nét liền;
\item Các đường bị khuất (gồm $1$ cạnh đáy và $1$ cạnh bên khuất phía sau) phải được vẽ bằng nét đứt khúc.
\end{itemize}
Do đó, hình ở phương án B là hình biểu diễn đúng.\\
(Hình A vẽ sai quy tắc vì toàn bộ nét đều là nét liền; hình C biểu diễn hình chóp tam giác/tứ diện; hình D là hình phẳng biểu diễn tứ giác cùng hai đường chéo).\\
\textbf{Chọn B.}
\end{tcolorbox}

\noindent\textbf{Câu 15. Lời giải chi tiết}\par
\begin{tcolorbox}[enhanced,breakable,colback=solutionbg,colframe=secondary!50!black,frame hidden,borderline west={3pt}{0pt}{secondary!50!black}]
Ta có:
\begin{itemize}[leftmargin=1.5em]
\item Điểm $S$ là điểm chung thứ nhất của hai mặt phẳng $(SAB)$ và $(SCD)$: $S \in (SAB) \cap (SCD)$.
\item Trong mặt phẳng đáy $(ABCD)$, điểm $E = AB \cap CD$. Vì $E \in AB \subset (SAB)$ và $E \in CD \subset (SCD)$ nên $E$ là điểm chung thứ hai: $E \in (SAB) \cap (SCD)$.
\end{itemize}
Do đó, giao tuyến của hai mặt phẳng $(SAB)$ và $(SCD)$ là đường thẳng $SE$.\\
\textbf{Chọn C.}
\end{tcolorbox}

\noindent\textbf{Câu 16. Lời giải chi tiết}\par
\begin{tcolorbox}[enhanced,breakable,colback=solutionbg,colframe=secondary!50!black,frame hidden,borderline west={3pt}{0pt}{secondary!50!black}]
Trong mặt phẳng $(SAC)$, gọi $P$ là giao điểm của hai đường thẳng $AM$ và $SO$. Khi đó:
$$\begin{cases} P \in AM \\ P \in SO \subset (SBD) \Rightarrow P \in (SBD) \end{cases} \Rightarrow P = AM \cap (SBD).$$
Vậy giao điểm của đường thẳng $AM$ với mặt phẳng $(SBD)$ là điểm $P$.\\
\textbf{Chọn D.}
\end{tcolorbox}

\noindent\textbf{Câu 17. Lời giải chi tiết}\par
\begin{tcolorbox}[enhanced,breakable,colback=solutionbg,colframe=secondary!50!black,frame hidden,borderline west={3pt}{0pt}{secondary!50!black}]
Xét hàm số lượng giác $y = f(x) = \tan x$:
\begin{itemize}[leftmargin=1.5em]
\item a) \textbf{Đúng.} Hàm số $y = \tan x = \dfrac{\sin x}{\cos x}$ xác định khi và chỉ khi:
$$\cos x \ne 0 \Leftrightarrow x \ne \dfrac{\pi}{2} + k\pi\quad (k\in\mathbb{Z}).$$
Vậy tập xác định của hàm số là $D = \mathbb{R}\setminus\left\{\dfrac{\pi}{2} + k\pi \;\middle|\; k\in\mathbb{Z}\right\}$.

\item b) \textbf{Đúng.} Thay $x_0 = \dfrac{\pi}{4}$ vào hàm số:
$$f\left(\dfrac{\pi}{4}\right) = \tan\dfrac{\pi}{4} = 1.$$

\item c) \textbf{Sai.} Hàm số lượng giác $y = \tan x$ là hàm số tuần hoàn với chu kì cơ sở là $T = \pi$, không phải $T = 2\pi$ (chu kì $2\pi$ là của các hàm số $\sin x$ và $\cos x$).

\item d) \textbf{Đúng.} Phương trình $f(x) = \sqrt{3} \Leftrightarrow \tan x = \sqrt{3} = \tan\dfrac{\pi}{3} \Leftrightarrow x = \dfrac{\pi}{3} + k\pi\quad (k\in\mathbb{Z})$.
Xét trên khoảng $\left(-\dfrac{\pi}{2}; \dfrac{\pi}{2}\right)$:
$$-\dfrac{\pi}{2} < \dfrac{\pi}{3} + k\pi < \dfrac{\pi}{2} \Leftrightarrow -\dfrac{1}{2} - \dfrac{1}{3} < k < \dfrac{1}{2} - \dfrac{1}{3} \Leftrightarrow -\dfrac{5}{6} < k < \dfrac{1}{6}.$$
Vì $k \in \mathbb{Z}$ nên chỉ có duy nhất $k = 0$, tương ứng với nghiệm duy nhất $x = \dfrac{\pi}{3} \in \left(-\dfrac{\pi}{2}; \dfrac{\pi}{2}\right)$.
\end{itemize}
\end{tcolorbox}

\noindent\textbf{Câu 18. Lời giải chi tiết}\par
\begin{tcolorbox}[enhanced,breakable,colback=solutionbg,colframe=secondary!50!black,frame hidden,borderline west={3pt}{0pt}{secondary!50!black}]
Số vốn ban đầu giải ngân ở quý 1 là $u_1 = 16$ tỷ đồng. Tỷ lệ tăng trưởng số vốn ở mỗi quý tiếp theo là $q = 1{,}5$.
\begin{itemize}[leftmargin=1.5em]
\item a) \textbf{Đúng.} Vì số vốn bơm vào dự án ở mỗi quý tiếp theo luôn gấp $1{,}5$ lần số vốn quý liền trước: $u_{n+1} = 1{,}5 \cdot u_n$ ($\forall n \ge 1$), nên dãy $(u_n)$ là một cấp số nhân có số hạng đầu $u_1 = 16$ tỷ đồng và công bội $q = 1{,}5$.

\item b) \textbf{Sai.} Theo công thức số hạng tổng quát của cấp số nhân:
$$u_n = u_1 \cdot q^{n-1} = 16 \cdot (1{,}5)^{n-1}\text{ (tỷ đồng)},\quad (\forall n \ge 1).$$
Khẳng định đưa ra công thức $u_n = 16 \cdot (1{,}5)^n$ là sai (công thức này cho $u_1 = 24$, không đúng với số vốn quý 1 là $16$).

\item c) \textbf{Đúng.}
\begin{itemize}
\item Số tiền giải ngân của riêng quý 4 là:
$$u_4 = 16 \cdot (1{,}5)^{4-1} = 16 \cdot (1{,}5)^3 = 16 \cdot \dfrac{27}{8} = 54\text{ (tỷ đồng)}.$$
\item Một năm gồm 4 quý, tổng số vốn đã giải ngân sau 1 năm đầu tiên là:
$$S_4 = u_1 + u_2 + u_3 + u_4 = 16 + 16 \cdot 1{,}5 + 16 \cdot (1{,}5)^2 + 54 = 16 + 24 + 36 + 54 = 130\text{ (tỷ đồng)}.$$
\end{itemize}

\item d) \textbf{Sai.} Tổng số vốn giải ngân lũy kế sau $n$ quý được tính bằng công thức tổng $n$ số hạng đầu của cấp số nhân:
$$S_n = \dfrac{u_1(q^n - 1)}{q - 1} = \dfrac{16[(1{,}5)^n - 1]}{1{,}5 - 1} = 32[(1{,}5)^n - 1]\text{ (tỷ đồng)}.$$
Sau đúng $6$ quý giải ngân ($n = 6$):
$$S_6 = 32[(1{,}5)^6 - 1] = 32\left[\left(\dfrac{3}{2}\right)^6 - 1\right] = 32\left[\dfrac{729}{64} - 1\right] = 32 \cdot \dfrac{665}{64} = \dfrac{665}{2} = 332{,}5\text{ (tỷ đồng)}.$$
Vì $332{,}5\text{ tỷ đồng} < 500\text{ tỷ đồng}$ nên sau 6 quý, tổng vốn giải ngân vẫn chưa vượt qua ngân sách dự kiến $500$ tỷ đồng.\\
(Thực tế, phải đến hết quý thứ 7: $S_7 = S_6 + u_7 = 332{,}5 + 16 \cdot (1{,}5)^6 = 332{,}5 + 182{,}25 = 514{,}75\text{ tỷ đồng} > 500\text{ tỷ đồng}$ thì tổng vốn mới vượt mức ngân sách). Do đó khẳng định d là sai.
\end{itemize}
\end{tcolorbox}

\noindent\textbf{Câu 19. Lời giải chi tiết}\par
\begin{tcolorbox}[enhanced,breakable,colback=solutionbg,colframe=secondary!50!black,frame hidden,borderline west={3pt}{0pt}{secondary!50!black}]
Dựa vào hình vẽ và các tính chất hình học không gian:
\begin{itemize}[leftmargin=1.5em]
\item a) \textbf{Đúng.} Điểm $G$ là trọng tâm của tam giác $SCD$ nên $G$ nằm trên đường trung tuyến $SM$ của tam giác $SCD$. Vì $SM \subset (SCD)$ nên $G \in (SCD)$.

\item b) \textbf{Đúng.} Hai mặt phẳng $(SCD)$ và $(SBD)$ có điểm chung thứ nhất là $S$ và điểm chung thứ hai là $D$. Do đó giao tuyến của $(SCD)$ và $(SBD)$ chính là đường thẳng $SD$.

\item c) \textbf{Đúng.} Đáy $ABCD$ là hình bình hành tâm $O$ nên $O$ là trung điểm của đường chéo $BD$. Mặt khác, $M$ là trung điểm của cạnh $CD$. Xét trong tam giác $BCD$, đoạn thẳng $OM$ nối hai trung điểm của hai cạnh $BD$ và $CD$ nên $OM$ là đường trung bình của tam giác $BCD$. Suy ra:
$$OM \parallel BC\quad \left(\text{và } OM = \dfrac{1}{2}BC\right).$$

\item d) \textbf{Sai.}
Xét tam giác $SMB$:
\begin{itemize}
\item Điểm $G$ là trọng tâm của tam giác $SCD$ trên đường trung tuyến $SM$ nên: $\dfrac{MG}{MS} = \dfrac{1}{3}$.
\item Điểm $E$ là trọng tâm của tam giác $BCD$ trên đường trung tuyến $BM$ nên: $\dfrac{ME}{MB} = \dfrac{1}{3}$.
\end{itemize}
Áp dụng định lí Thalès đảo trong tam giác $SMB$, ta có:
$$\dfrac{MG}{MS} = \dfrac{ME}{MB} = \dfrac{1}{3} \Rightarrow GE \parallel SB.$$
Vì đường thẳng $GE$ song song với đường thẳng $SB$ mà $SB \subset (SAB)$ và $GE \not\subset (SAB)$, nên đường thẳng $GE$ song song với mặt phẳng $(SAB)$:
$$GE \parallel (SAB).$$
Do $GE \parallel (SAB)$ nên đường thẳng $GE$ không thể cắt mặt phẳng $(SAB)$, khẳng định đường thẳng $GE$ cắt $(SAB)$ tại điểm $P \in SB$ là sai.
\end{itemize}
\end{tcolorbox}

\noindent\textbf{Câu 20. Lời giải chi tiết}\par
\begin{tcolorbox}[enhanced,breakable,colback=solutionbg,colframe=secondary!50!black,frame hidden,borderline west={3pt}{0pt}{secondary!50!black}]
Giải phương trình lượng giác: $\cos\left(2x + \dfrac{\pi}{4}\right) = -\dfrac{\sqrt{2}}{2}$.
\begin{itemize}[leftmargin=1.5em]
\item \textbf{Biến đổi phương trình về dạng chuẩn và viết họ nghiệm (0,5 điểm):}\\
Ta có $-\dfrac{\sqrt{2}}{2} = \cos\left(\dfrac{3\pi}{4}\right)$, phương trình tương đương:
$$\cos\left(2x + \dfrac{\pi}{4}\right) = \cos\left(\dfrac{3\pi}{4}\right) \Leftrightarrow \left[\begin{array}{l}
2x + \dfrac{\pi}{4} = \dfrac{3\pi}{4} + k2\pi \\[6pt]
2x + \dfrac{\pi}{4} = -\dfrac{3\pi}{4} + k2\pi
\end{array}\right. \quad (k \in \mathbb{Z}).$$

\item \textbf{Giải chi tiết tìm $x$ (0,5 điểm):}
$$\Leftrightarrow \left[\begin{array}{l}
2x = \dfrac{3\pi}{4} - \dfrac{\pi}{4} + k2\pi \\[6pt]
2x = -\dfrac{3\pi}{4} - \dfrac{\pi}{4} + k2\pi
\end{array}\right. \Leftrightarrow \left[\begin{array}{l}
2x = \dfrac{\pi}{2} + k2\pi \\[6pt]
2x = -\pi + k2\pi
\end{array}\right. \Leftrightarrow \left[\begin{array}{l}
x = \dfrac{\pi}{4} + k\pi \\[6pt]
x = -\dfrac{\pi}{2} + k\pi
\end{array}\right. \quad (k \in \mathbb{Z}).$$

\item \textbf{Kết luận đầy đủ họ nghiệm kèm điều kiện (0,5 điểm):}\\
Vậy tập nghiệm của phương trình đã cho là:
$$S = \left\{\dfrac{\pi}{4} + k\pi;\; -\dfrac{\pi}{2} + k\pi \;\middle|\; k \in \mathbb{Z}\right\}.$$
\end{itemize}
\end{tcolorbox}

\noindent\textbf{Câu 21. Lời giải chi tiết}\par
\begin{tcolorbox}[enhanced,breakable,colback=solutionbg,colframe=secondary!50!black,frame hidden,borderline west={3pt}{0pt}{secondary!50!black}]
\textbf{a) Vẽ hình biểu diễn chính xác, phân biệt rõ nét đứt/liền (0,5 điểm):}
\begin{center}
\begin{tikzpicture}[scale=0.9, line join=round, line cap=round]
  % Toa do cac dinh hinh chóp day hinh binh hanh
  \coordinate (A) at (0,0);
  \coordinate (B) at (4.6,-0.6);
  \coordinate (D) at (1.6,1.8);
  \coordinate (C) at ($(B)+(D)-(A)$);
  \coordinate (S) at (2.2,4.8);
  
  % Giao diem O cua AC va BD
  \coordinate (O) at (intersection of A--C and B--D);
  % Diem M la trung diem SC
  \coordinate (M) at ($(S)!0.5!(C)$);
  % Giao diem I cua AM va SO trong mp (SAC)
  \coordinate (I) at (intersection of A--M and S--O);
  % Giao diem K cua BI va SD trong mp (SBD)
  \coordinate (K) at (intersection of B--I and S--D);
  
  % Net dut khuat
  \draw[dashed, thick] (A)--(D)--(C) (S)--(D) (A)--(C) (B)--(D) (S)--(O);
  \draw[dashed, thick, primary] (A)--(M) (B)--(I);
  \draw[dashed, thick, accent] (I)--(K) (M)--(K);
  
  % Net lien nhin thay
  \draw[thick] (S)--(A)--(B)--(C)--(S)--(B);
  \draw[thick, primary] (B)--(M);
  
  % Cac diem
  \fill (S) circle (2pt) node[above] {$S$};
  \fill (A) circle (2pt) node[below left] {$A$};
  \fill (B) circle (2pt) node[below right] {$B$};
  \fill (C) circle (2pt) node[right] {$C$};
  \fill (D) circle (2pt) node[above left] {$D$};
  \fill (O) circle (2pt) node[below] {$O$};
  \fill[primary] (M) circle (2.2pt) node[above right] {$M$};
  \fill[primary] (I) circle (2pt) node[below left] {$I$};
  \fill[accent] (K) circle (2.2pt) node[above left] {$K$};
\end{tikzpicture}
\end{center}

\textbf{b) Lập luận chọn mặt phẳng phụ và tìm giao tuyến trung gian (0,5 điểm):}
\begin{itemize}[leftmargin=1.5em]
\item Để xác định giao điểm của đường thẳng $SD$ với mặt phẳng $(ABM)$, ta chọn mặt phẳng phụ chứa đường thẳng $SD$ là $(SBD)$.
\item Tìm giao tuyến của mặt phẳng phụ $(SBD)$ và mặt phẳng $(ABM)$:\\
+ Ta thấy điểm $B$ thuộc cả hai mặt phẳng: $B \in (SBD)$ và $B \in (ABM)$ nên $B$ là điểm chung thứ nhất:
$$B \in (SBD) \cap (ABM).$$
+ Trong mặt phẳng $(SAC)$, gọi $I$ là giao điểm của hai đường thẳng $SO$ và $AM$:
$$I = SO \cap AM.$$
Vì $I \in AM \subset (ABM)$ và $I \in SO \subset (SBD)$ nên $I$ là điểm chung thứ hai:
$$I \in (SBD) \cap (ABM).$$
+ Do đó, giao tuyến của hai mặt phẳng $(SBD)$ và $(ABM)$ là đường thẳng $BI$:
$$(SBD) \cap (ABM) = BI.$$
\end{itemize}

\textbf{c) Xác định chính xác giao điểm và kết luận (0,5 điểm):}
\begin{itemize}[leftmargin=1.5em]
\item Trong mặt phẳng $(SBD)$, đường thẳng $BI$ cắt đường thẳng $SD$ tại điểm $K$:
$$K = BI \cap SD.$$
\item Khi đó:
$$\begin{cases} K \in SD \\ K \in BI \subset (ABM) \Rightarrow K \in (ABM) \end{cases} \Rightarrow K = SD \cap (ABM).$$
Vậy giao điểm của đường thẳng $SD$ với mặt phẳng $(ABM)$ là điểm $K$.\\
\textit{(Ghi chú thêm: Trong $\triangle SAC$, $SO$ và $AM$ là hai đường trung tuyến cắt nhau tại $I$ nên $I$ là trọng tâm $\triangle SAC \Rightarrow \dfrac{SI}{SO} = \dfrac{2}{3}$. Do đó trong $\triangle SBD$, $I$ cũng là trọng tâm tam giác, suy ra $BI$ là đường trung tuyến của $\triangle SBD$ ứng với đỉnh $B$, do đó $K$ chính là trung điểm của cạnh bên $SD$).}
\end{itemize}
\end{tcolorbox}

\end{document}
